\paragraph{A-CS (Open; common admissible sources over $\mathbb Z$).}
For arbitrary two admissible strict generalized PD de Rham theories
$A,B$ in the category specified in Part A, we have not constructed
an admissible $D$ with strict coefficient-preserving maps
$D\to A$ and $D\to B$, nor exhibited two objects for which such
a source cannot exist.  This includes the version with support
$q\leq w$ and the version with termwise flat integral groups.
We keep these versions in scope; the examples and refinements in
this paper explicitly record when both extra properties hold.

The strongest proved common-source reduction is A-S09-29.
For each fixed $d\geq0$, source and targets may be taken in the
proper full subcategory $\mathcal E^{[d]}_\omega\subset\mathcal C^{[d]}$.
Every homogeneous group is countable, with no finite-generation
or computability assertion.
Its first skeleton is the specified $P(N)$ with one disk in every
weight at least $N$; its coefficient comparison holds through $d$
in every weight, and its rows through $d+2$ are exact from $N$.
The two targets can have the same $N$, while the internal source
has its own cutoff. All replacement arrows and both directions
are explicit. This internal common-source assertion remains open.
The terminal-localization gate is now A-S09-30, with the actual
one-copy, two-copy and three-copy comparison diagrams and their
finite-word localization comparison.
Each reduction concerns a fixed finite $d$; no single object exact
in every degree is produced by these comparisons.

The preceding A-006 reduction uses $\mathcal C$
be the proper full subcategory defined there by the first-skeleton
property $P(E)$.  The actual functor $R:\mathcal A\to\mathcal C$
and natural comparison $jR\Rightarrow\mathrm{Id}$ prove that
common sources for arbitrary admissible objects are equivalent
to common sources internal to $\mathcal C$ for every pair of
its objects.  This assertion for $\mathcal C$ is undecided.
A-003 additionally identifies it with connectivity by finite
zigzags.  No such zigzag between arbitrary objects of
$\mathcal C$ has been constructed.  This is a genuine restriction
of both the objects and the allowed source, with every comparison
map and both directions proved in A-006.
The sufficient arithmetic completion construction of SA08 C30/C34
still lacks local repairs with the required compatible target maps.
AT12 refutes the assertion that \emph{every prescribed} C21 partial
diagram has repairs preserving forms zero and one at the indicated
primes.  The example retains the explicit maps $G\to A\to G$.
Consequently it does not refute common-source existence or eliminate
the possibility of choosing a different partial diagram.

The smallest remaining construction step in the approaches pursued
here is an unscaled relative PD attachment with simultaneous maps to
both arbitrary targets, controlled vertical cohomology, and uniform
eventual horizontal bounds.  A-L05 and A-S09-13 extend coefficient
or degree-zero-factor scaled disks through the full bar construction.
Their formulas supply precisely those prescribed maps; they do not
factor an arbitrary boundary.  AT10 gives an actual tail lattice and
its split-parameter pushouts, with the exact factorial coefficients
in (AT.26).  A-H01 shows that the coefficient-tail multiples of
horizontal classes supplied by those particular evaluations already
vanish under the original row axiom.  Thus these evaluations have
not filled a general unscaled class.

An alternative weight induction has several proved steps:
A-S09-14 removes torsion $H^1$ at a single weight by a canonical
based saturation and extends both target maps.  Its PD-algebra
version requires the stated torsion-free terms.  It leaves free
$H^1$, higher vertical groups, later-weight products and uniform
admissibility unproved.  Those are retained tasks, rather than
conclusions of the single-weight result.

AT13--15 now give actual coefficient-resolution and two-cell
attachments with extending target maps and their exact ranges.
A-S09-19 chooses the initial C21 source with $P_0=S$ and with
$P_1$ an integral coefficient resolution.  AT15 preserves both
properties for the specified relative two-simplex repairs in
form degree two and in every form degree at least four.  Its
degree-three case has the exact residue (AT.65).  Horizontal
restoration after each permitted repair is constructed there.
These individual steps do not yet prove a simultaneous iteration
with all required eventual vertical bounds.

A-S09-18 identifies the initial degree-four group naturally;
A-S09-20 computes the full low cohomology of its independent
free-parameter repair in every weight.  A-S09-21 gives a tail
algebra with the stated low cohomology, together with the precise
two $S$-linear equations still needed for its target maps.
AT16--17 construct the coefficient-resolution attachments with
actual target maps, including the prescribed edge boundary.
A-SAT01 supplies the iterated degree-one saturation and extends
all specified torsion-free target maps.  The direct comparison
with the two-term tail has now been refuted as a low-cohomology
equivalence: A-L11 gives an order-two degree-one class, and A-L12
and A-S09-22 give order-two classes in degrees two and three that
survive every based saturation stage.  AT18 gives the exact
coefficient residue for attaching the next degree-three repair,
and proves its stated invariance with the extra low-weight
period hypothesis explicitly retained.  AT19 supplies actual zero-residue lifts and repairs for the
genuine quadratic sphere disk under its explicit bounds.
A-L13 constructs the variable-boundary pair repair exactly when
the two residues agree.  A general simultaneous vanishing or
compatible repair for arbitrary subsequent source data is unproved.
A-S09-23 also refutes a second-skeleton extension using only
one-forms generated by zero-forms and their differentials, for
$N\geq2$.  The free first-skeleton reduction itself is unaffected.
No completed admissible common source follows from these steps.

\paragraph{A-TL (Open; terminal localization over $\mathbb Z$).}
Let $\mathcal A$ denote this category and invert all its morphisms.
A-006 supplies explicit mutually inverse equivalences between
this groupoid and $\mathcal C[\operatorname{Mor}(\mathcal C)^{-1}]$.
A-S09-24 transfers this question to $\mathcal C^{[d]}$,
A-S09-29 to its countable full subcategory, and A-S09-30 further
to $\mathcal F_{3,\omega}^{[d]}$.
The last category has first skeleton $P_r(N)$ with
$r\in\{1,2,3\}$, global coefficient comparison through $d$,
and rows through $d+2$ exact from $N$. Its localization is
equivalent to the original one by the exact copy-inclusion,
fold and three-copy composition diagrams. This is the current
exact terminal-localization gate for every fixed finite $d$.
A-S09-28 proves that the smaller single-copy localization has
a nonidentity endomorphism whose class is killed by the
explicit doubled-skeleton diagram. That calculation proves
the failure of the particular single-copy localization
comparison; it gives no nonidentity class in the original
localized category.
We have not proved that its resulting groupoid is equivalent to the
one-object groupoid with trivial automorphism group.  A-003 proves
that such a conclusion supplies common sources, since an isomorphism
in the groupoid is represented by a finite zigzag.  We have not
proved the converse: triviality of all zigzag loops requires further
equalization or homotopy data.  The pair-target and homotopy-equalizer
repairs in SA08 C32 are still unconstructed in the required integral
form.  None of the explicitly split examples in this paper has
provided a nontrivial invariant of such loops.

\paragraph{A-SA09 (Open application; reconstruction theorems proved).}
A-S09-01, A-S09-02 and A-S09-03 independently establish the stated
SA09 A2.2, A3.1 and A3.2, with their exact categories and essential
images.  The remaining task is an actual construction from those
structures that supplies a common admissible source, contracts the
localized category, or gives an invariant separating its objects.
A-S09-04 gives a split admissible morphism on which the raw row-test
assignment changes.  It refutes that assignment as the proposed
invariant.  It does not exclude a different derived construction.
A-S09-06--08 refute the three particular strict rectifications
stated there.  A-S09-09--12 calculate unbounded cotangent support
for each fixed nontrivial coefficient truncation, with all comparison
maps and source hypotheses included in their proofs.  No bounded
resolution or strict comparison for arbitrary admissible targets is
inferred from those calculations.

All three open questions remain part of this incomplete investigation.
The results register and progress ledger retain them separately.
